Das Polynom $m(X)=X^2-X-1\in\mathbb{Q}[X]$ ist irreduzibel und definiert damit
ein algebraisches Element $\varphi$ über $\mathbb{Q}$.
Führen Sie die folgenden Berechnungen im Körper $\mathbb{Q}(\varphi)$
durch.
\begin{teilaufgaben}
\item Berechnen Sie $(2+\varphi)(3-2\varphi)$.
\item Berechnen Sie $\varphi^{-1}$.
\item Berechnen Sie $(3-2\varphi)^{-1}$.
\item Berechnen Sie $(2+\varphi)/(3-2\varphi)$.
\item Löst man die quadratische Gleichung, findet man zwei Lösungen
\[
\varphi_\pm
=
\frac{1\pm\!\sqrt{5}}2,
\]
die beide mögliche Werte für $\varphi$ sind und daher die gleichen
algebraischen Eigenschaften haben.
Prüfen Sie nach, dass $1/\varphi_\pm = -\varphi_\mp$.
\end{teilaufgaben}

\begin{hinweis}
Das Element $\varphi_+=1.618\dots $ ist das Verhältnis des goldenen Schnitts.
\end{hinweis}

\begin{loesung}
\begin{teilaufgaben}
\item
Ausmultiplizieren ergibt
\begin{align*}
(2+\varphi)(3-2\varphi)
&=
6-4\varphi+3\varphi-2\varphi^2
=
6-\varphi-2\varphi^2.
\intertext{Das Quadrat von $\varphi$ kann dank der definierenden Gleichung
für $\varphi$ durch $\varphi^2=\varphi+1$ ersetzt werden, was}
&=
6-\varphi-2(\varphi+1)
=5-3\varphi
\end{align*}
ergibt.

\item
Wir multiplizieren die definierende Gleichung $\varphi^2-\varphi-1=0$
mit $\varphi^{-1}$ und lösen nach $\varphi^{-1}$ auf:
\[
(\varphi^2-\varphi-1)\varphi^{-1}=0
\quad\Rightarrow\quad
\varphi-1-\varphi^{-1}=0
\quad\Rightarrow\quad
\varphi^{-1}=\varphi-1.
\]

\item 
Gesucht ist $a+b\varphi$ mit
\begin{align*}
(3-2\varphi)(a+b\varphi)&=1.
\intertext{Ausmultiplizieren ergibt}
3a +(3b-2a)\varphi -2b\varphi^2&=1,
\intertext{worin wir wieder $\varphi^2=\varphi+1$ ersetzen können und}
3a +(3b-2a)\varphi -2b(\varphi+1)&=1\\
(3a-2b) +(b-2a)\varphi &=1
\end{align*}
erhalten.
Daraus lesen wir die Gleichungen
\[
\left.
\renewcommand{\arraycolsep}{2pt}
\begin{array}{rcrcr}
 3a &-& 2b &=& 1 \\
-2a &+&  b &=& 0
\end{array}
\quad
\right\}
\qquad\text{mit den Lösungen}\qquad
\left\{\quad
\begin{aligned}
a&=-1
\\
b&=-2
\end{aligned}
\right.
\]
ab.
Es folgt $(3-2\varphi)^{-1}=-1-2\varphi$.
Wir kontrollieren das Resultat durch Multiplikation
\[
(3-2\varphi)(-1-2\varphi)
=
-3+2\varphi-6\varphi + 4\varphi^2
=
-3-4\varphi+4(\varphi+1)
=
1.
\]

\item
Wir multiplizieren mit der in Teilaufgabe c) gefundenen multiplikativen
Inversen
\begin{align*}
\frac{2+\varphi}{3-2\varphi}
&=
(2+\varphi)(-1-2\varphi)
=
-2-\varphi-4\varphi-2\varphi^2
=
-2-5\varphi-2(\varphi+1)
=
-4-7\varphi.
\end{align*}

\item
Durch Rationalisieren finden wir
\begin{align*}
\varphi_\pm^{-1}
&=
\frac{1}{\varphi_\pm}
=
\frac{2}{1\pm\!\sqrt{5}}
=
\frac{2}{1\pm\!\sqrt{5}}
\frac{1\mp\!\sqrt{5}}{1\mp\!\sqrt{5}}
=
\frac{2(1\mp\!\sqrt{5})}{1-5}
=
-\frac{1\mp\!\sqrt{5}}{2}
=
-\varphi_\mp.
\qedhere
\end{align*}
\end{teilaufgaben}
\end{loesung}
